\documentclass[11pt,a4paper]{article}
\usepackage[margin=23mm]{geometry}
\usepackage{amsmath,amssymb,booktabs,array}
\usepackage[hidelinks]{hyperref}
\setlength{\parindent}{0pt}
\setlength{\parskip}{5pt}
\emergencystretch=3em
\raggedright
\begin{document}
\begin{center}
{\Large BP-R Stage 3: source laws and merge choices}\\[4pt]
6 October 2026 --- supplied originals read; S3-G2-N open
\end{center}
\textbf{Scope.} Source transcriptions and static proposed-node evaluations.
Neither node is adopted. No simulator, traffic experiment or gradient validation.

\section*{1. What the supplied journal actually specifies}
Laval and Daganzo (2006), Section 2.1, equations 6--13, specifies continuous
lane traffic with a triangular fundamental diagram. In this unit-consistent
transcription, $k_\ell$ is donor density (veh/m), $u$ free-flow speed (m/s),
$w>0$ backward wave-speed magnitude (m/s), $\kappa$ jam density (veh/m),
$Q$ lane capacity (veh/s), and $\Delta t$ a time interval (s). Sending mass:
\begin{equation}
S_\ell=\Delta t\min\{u k_\ell,Q\}.
\end{equation}
For allowed target lane $b$, lane-choice speeds $\bar v$ (m/s) and relaxation
time $\tau$ (s), desired lane-change rate and transfer masses are
\begin{align}
\pi_{\ell b}&=\frac{(\bar v_b-\bar v_\ell)_+}{u\tau},&
d_{\ell b}&=\pi_{\ell b}\Delta t\,S_\ell,\quad b\ne\ell,\\
d_{\ell\ell}&=\left(1-\sum_{b\ne\ell}\pi_{\ell b}\Delta t\right)S_\ell.
\end{align}
$(z)_+=\max(0,z)$. The step must give
$\sum_{b\ne\ell}\pi_{\ell b}\Delta t\le1$; outgoing demands sum to $S_\ell$.
Source lane changes advance from cell $(i,\ell)$ to $(i+1,b)$;
they are not counted twice as separate lateral/longitudinal moves.

For target cell/lane $b$, collect all compatible incoming demands $d_{jb}$
and form receiving mass and actual transfers:
\begin{equation}
R_b=\Delta t\min\{w(\kappa-k_b),Q\},\qquad
\gamma_b=\min\left\{1,\frac{R_b}{\sum_jd_{jb}}\right\},\qquad
f_{jb}=\gamma_b d_{jb}.
\end{equation}
If demand is zero, all transfers are zero; no ratio is evaluated. This
transcribes the journal's local IT allocation in mass units. Transfers respect
receiving and desired donor demand. Pair each debit and credit:
\begin{equation}
n_c^{+}=n_c+\sum_j f_{jc}-\sum_b f_{cb}.
\end{equation}
Boundary admissions/exits are separate. Internal transfers cancel in the sum;
event counts never round $n_c$.

\textbf{Speed and receiving scope.} Lane-choice speed uses the journal's
downstream look-ahead/preceding-time space-mean convention, not arbitrary point
FD inversion. Equation (4) is the unobstructed local recipe. Active moving
constraints must use the coupled source method; do not reuse unobstructed
supply and add a guessed acceleration penalty.

\newpage
\section*{2. The motion law is now source-backed}
Appendix A, equations A.1--A.4, uses disturbance position $\phi(t)$ (m),
realized speed $v(t)$ (m/s), desired kinetic speed $v_{\rm des}(t)$, and
immediately downstream KW-stream speed $v_{\rm dwn}(t)$:
\begin{align}
v(t)&=\min\{v_{\rm dwn}(t),v_{\rm des}(t)\},\\
v_{\rm des}(t+\Delta t)&=v(t)+a(v(t),\phi(t))\Delta t,\\
\phi(t+\Delta t)&=\phi(t)+v(t)\Delta t,\\
a(v,x)&=a_0(1-v/v_{\max})-gG(x).
\end{align}
$a_0$ is zero-speed acceleration (m/s$^2$), $v_{\max}$ vehicle free maximum
speed (m/s), $g$ gravity (m/s$^2$), and $G(x)$ dimensionless grade. Traffic
constrains realized speed; $v_{\max}$ need not equal $u$. Numerical examples
in the journal and merge manuscript have different parameters; none is
silently combined or calibrated.

$v_{\rm dwn}$ is a numerical estimate over the ensuing interval on the
occupied lane, computed from the KW problem affected by the constraint.
It is not independently assigned from the upstream command. Zero passing
\emph{relative to the moving disturbance} permits positive stationary-detector flow.

For cumulative actual transfer mass $H(t)$, the deterministic option creates
descriptors when $\lfloor H(t)\rfloor$ increases. Only descriptors are discrete;
the conserved transfer in (5) remains continuous, including fractional mass.
Traffic is already counted in the field. Descriptors are not extra TST
population. A complete task must declare birth speed, event phase and insertion
timing. No random sampling or coincidence jitter is approved.

\textbf{Numerical lineage.} Section 2.2 cites the 2005 method that converges
in flows; Appendix A cites the separate numerical-treatment paper. Available
2003 originals distinguish cumulative-count approximation from continuous-space
composite-Riemann coupling. The latter uses wave-interaction and crossing step
conditions. This supports a conservative route, not a theorem for every
endogenous event/control derivative.

\textbf{Second supplied paper.} Daganzo, Lin and Del Castillo (1997),
Sections 2--3, proves IT equivalence for two classes with different permissions
to two lane sets. Its priority vehicles are an access class. Section 4
illustrates an exit-ramp/diverge; Section 5 leaves general junction extension
open. It does not uniquely determine the merge manuscript's 40\% adjustment.

\textbf{Preserved task.} Control caps admission to the finite KW ramp approach.
Merge entry follows later demand and receiving; it is not also capped by the
contemporaneous command. All waiting, ramp and mainline inventory is in TST.
There is no added recovery-cost term.

\newpage
\section*{3. Remaining choice: mandatory point merge}
$D_r,D_m\ge0$ are ramp and non-ramp incoming demand rates (veh/s);
$R\ge0$ is shared target receiving (veh/s), including existing moving-constraint
restrictions. Non-ramp demand includes all compatible through/lateral inflows.
No other movement may spend $R$ again. Actual rates are $f_r,f_m$.

\textbf{A: fixed junction priority.} A fixed ramp share $p_r\in(0,1)$ is
common to all controls/models. Daganzo's 1994 original, Section 3.2, supplies
the two-approach construction, transcribed as
\begin{equation}
(f_r,f_m)=
\begin{cases}
(D_r,D_m),&D_r+D_m\le R,\\
\left(\min\{D_r,\max(p_rR,R-D_m)\},\,R-f_r\right),
&D_r+D_m>R.
\end{cases}
\end{equation}
Unused shares are reallocated. Grouping local lane demands and applying this
node to the selected reference is a proposed adaptation, not exact reproduction
of the unread gridlock chapter. No $p_r$ value is inferred from 40\%.

\textbf{B: demand-proportional allocation.} Extend local IT to the mandatory node:
\begin{equation}
(f_r,f_m)=
\begin{cases}
(D_r,D_m),&D_r+D_m\le R,\\
\displaystyle\frac{R}{D_r+D_m}(D_r,D_m),&D_r+D_m>R.
\end{cases}
\end{equation}
No fixed share is introduced. Contested fractions change with demands. This
extension is not a general-junction theorem from the 1997 paper.

Both give $0\le f_j\le D_j$ and
$f_r+f_m=\min(D_r+D_m,R)$. In (10)'s overloaded case,
$R-D_m\le f_r\le D_r$ when the lower bound is positive, and
$0\le f_r\le R$; hence $0\le R-f_r\le D_m$. In (11), the common multiplier
is between zero and one. These are static algebraic checks, not traffic runs.

\begin{center}
\begin{tabular}{rrrrr}
\toprule
$R$ & $D_m$ & $D_r$ & A: $f_m,f_r$ & B: $f_m,f_r$\\
\midrule
100 & 90 & 90 & 70,\ 30 & 50,\ 50\\
100 & 20 & 90 & 20,\ 80 & 18.18,\ 81.82\\
100 & 90 & 10 & 90,\ 10 & 90,\ 10\\
100 & 0 & 90 & 0,\ 90 & 0,\ 90\\
0 & 90 & 90 & 0,\ 0 & 0,\ 0\\
\bottomrule
\end{tabular}
\end{center}
Illustrative consistent rate units; A uses $p_r=0.30$ only for these examples.
Proportional allocation can leave some mainline demand unserved although it
is below A's fixed share. Both use all available receiving.

\textbf{Recommendation and limits.} A has direct merge-specific primary support.
B avoids a share parameter but assumes demand-based competition. Either needs
adoption at S3-G2-N. Then qualify joined event/receiving timing, clearance,
inactive equivalence, strong independent closures and actual rate leverage.
No full-TST benefit, fine-state necessity, useful gradient or Stage 3 closure
is established here.
\end{document}

